\chapter{Resultados e Discussões}

Este capitulo apresenta os resultados finais dos quatro modelos avaliados para deteccao de intrusoes em rede CAN: Regressao Logistica, SVM, MLP e XGBoost. A base canonica utilizada corresponde aos artefatos \texttt{run\_summary.json} e \texttt{classification\_report.txt} dos diretorios \texttt{*\_final\_clean}.

\section{Comparacao Principal entre Modelos}

A Tabela \ref{tab:comparacao_principal} consolida as metricas globais de comparacao no conjunto de teste.

\begin{table}[H]
\centering
\caption{Comparacao principal entre os modelos (resultados finais limpos)}
\begin{tabular}{|l|c|c|c|c|}
\hline
\textbf{Modelo} & \textbf{Accuracy} & \textbf{Precision Macro} & \textbf{Recall Macro} & \textbf{F1 Macro} \\
\hline
Regressao Logistica & 0.97252 & 0.97587 & 0.96940 & 0.97235 \\
SVM (LinearSVC) & 0.94992 & 0.94689 & 0.95387 & 0.95000 \\
MLP & 0.99908 & 0.99937 & 0.99873 & 0.99905 \\
XGBoost & 0.99998 & 0.99999 & 0.99998 & 0.99998 \\
\hline
\end{tabular}
\label{tab:comparacao_principal}
\caption*{Fonte: elaborada pelo autor com base em \texttt{classification\_report.txt} dos experimentos finais.}
\end{table}

Os resultados evidenciam superioridade de XGBoost no cenario avaliado, com desempenho quase perfeito nas metricas globais. O MLP apresentou resultado muito proximo e tambem altamente competitivo. Regressao Logistica manteve desempenho forte para um modelo linear, enquanto o SVM ficou abaixo dos demais nas metricas globais.

\section{Analise por Classe}

Além das metricas agregadas, a avaliacao por classe mostra que \texttt{GEAR} e \texttt{RPM} foram classes com separacao quase perfeita na maior parte dos modelos. Em contrapartida, \texttt{Fuzzy} foi a classe mais desafiadora, concentrando as menores taxas de recall entre os quatro algoritmos.

No caso da Regressao Logistica, o recall de Fuzzy foi 0.89292. No SVM, esse valor foi 0.88677, representando a pior taxa de recall entre os modelos. Para o MLP, o recall de Fuzzy subiu para 0.99368, e no XGBoost atingiu 0.99990.

Esse comportamento reforca que a dificuldade de classificacao nao esta distribuida de forma uniforme entre classes. Assim, a discussao de desempenho por classe e indispensavel para uma conclusao robusta sobre IDS multiclasse em CAN.

\section{Discussao dos Resultados}

Os resultados permitem destacar tres pontos principais.

Primeiro, ha um ganho expressivo ao migrar de modelos lineares para modelos nao lineares e de maior capacidade representacional, como MLP e XGBoost.

Segundo, mesmo quando a acuracia global e elevada, a avaliacao por classe continua necessaria para revelar classes mais sensiveis, como Fuzzy.

Terceiro, no recorte experimental adotado, XGBoost apresentou o melhor compromisso global, seguido de perto por MLP, tornando-se a referencia principal para recomendacao no capitulo de conclusoes.

\section{Sintese do Capitulo}

Com base na comparacao final, o ranking global de desempenho no cenario estudado foi:

\begin{enumerate}
    \item XGBoost;
    \item MLP;
    \item Regressao Logistica;
    \item SVM.
\end{enumerate}

Essa sintese responde ao objetivo comparativo do trabalho e fundamenta a recomendacao final do modelo mais adequado para IDS em rede CAN no contexto deste experimento.

